\section{Kapitel 3}
\subsection{Bibeltext}

        \spa{2}\hh{1}gleicherweise sollen die Frauen sich \\
            \spa{3}den eigenen Männern unterordnen, \\
            \spa{3}damit auch, \\
            \spa{3}wenn etliche dem Wort im Unglauben ungehorsam sind, \\
            \spa{3}sie durch [ihrer] Frauen Lebensführung ohne Wort gewonnen werden,\\
            \spa{3}\hh{2}wenn sie eure in Furcht geführte reine Lebensführung beobachtet haben,\\
            \spa{3}\hh{3}deren Schmuck nicht der äußere sei \\
                \spa{4}– Haarflechten und \\
                \spa{4}Goldumhängen \\
                \spa{4}oder Kleideranziehen –,\\
            \spa{3}\hh{4}sondern der verborgene Mensch des Herzens \\
                \spa{4}in dem Unverderblichen des sanften und stillen Geistes, \\
                \spa{4}der vor Gott von hohem Wert ist;\\
            \spa{3}\hh{5}denn so schmückten sich auch einst die heiligen Frauen, \\
                \spa{4}die auf Gott hofften \\
                \spa{4}und sich den eigenen Männern unterordneten,\\
                    \spa{5}\hh{6}wie Sara dem Abraham gehorchte \\
                    \spa{5}und ihn „Herr“ nannte, \\
                \spa{4}deren Kinder ihr wurdet als solche, \\
                \spa{4}die Gutes tun \\
                \spa{4}und nicht fürchten \\
                \spa{4}irgendeinen Schrecken; \\
        \spa{2}\hh{7}gleicherweise die Männer: \\
            \spa{3}Wohnt der Kenntnis entsprechend mit dem weiblichen Gefäß \\
            \spa{3}als dem schwächeren zusammen \\
            \spa{3}und erteilt ihm Ehre als die \\
                \spa{4}ihr auch Miterben seid der Gnade des Lebens, \\
                \spa{4}sodass eure Gebete nicht abgeschnitten werden; \\
        \spa{2}\hh{8}schließlich, \\
            \spa{3}[seid] alle gleichgesinnt, \\
                \spa{4}mitleidig, \\
                \spa{4}brüderlich liebend, \\
                \spa{4}herzlich und feinfühlig, \\
                \spa{4}freundlich gesinnt\\
                \spa{4}\hh{9}und vergeltet nicht Böses mit Bösem \\
                \spa{4}oder Schimpfwort mit Schimpfwort – \\
            \spa{3}im Gegenteil: \\
                \spa{4}Segnet, \\
                \spa{4}in dem Wissen, \\
                \spa{4}dass ihr hierzu gerufen wurdet, \\
                \spa{4}damit ihr Segen erbt; \\
            \spa{3}\hh{10}denn wer [das] Leben lieben und gute Tage sehen will, \\
                \spa{4}bringe seine Zunge dazu, \\
                \spa{4}vom Bösen zu lassen, \\
                \spa{4}und seine Lippen, \\
                \spa{4}dass sie nicht Trug reden. \\
            \spa{3}\hh{11}Er biege ab vom Bösen und tue Gutes. \\
            \spa{3}Er s\spa{4}uche Frieden und jage ihm nach –, \\
                \spa{4}\hh{12}weil „des Herrn Augen auf die Gerechten [gerichtet sind] \\
                \spa{4}und seine Ohren zu ihrem Flehen. \\
                \spa{4}Aber das Angesicht des Herrn ist auf die [gerichtet], \\
                \spa{4}die Böses tun.“ {Ps 34,16.17 n. d. gr. Üsg.} \\

        \spa{2}\hh{13}Und wer ist es, \\
            \spa{3}der euch schaden wird, \\
            \spa{3}wenn ihr Nachahmer des Guten werdet? \\
        \spa{2}\hh{14}Wenn ihr jedoch auch wegen Gerechtigkeit zu leiden habt – \\
            \spa{3}Selige [seid ihr]! \\
            \spa{3}Ihre Furcht fürchtet nicht. \\
            \spa{3}Lasst euch auch nicht in Unruhe versetzen. \\
                \spa{4}\hh{15}Den Herrn aber, \\
                \spa{4}Gott, heiligt in euren Herzen. \\

    \spa{1}Seid immer bereit zu einer Verteidigung vor jedem, \\
        \spa{2}der Rechenschaft fordert bezüglich der Hoffnung in euch, \\
        \spa{2}[und das] mit Sanftmut und Furcht! \\
        \spa{2}\hh{16}Und habt ein gutes Gewissen, \\
            \spa{3}damit die, \\
            \spa{3}die eure gute Lebensführung in Christus verunglimpfen, \\
            \spa{3}in dem Reden gegen euch als [vorgebliche] Übeltäter beschämt werden, \\
        \spa{2}\hh{17}denn es ist besser, \\
            \spa{3}wenn Gottes Wille es [so haben] will, \\
            \spa{3}dass ihr für Gutestun leidet als für Bösestun, \\
                \spa{4}\hh{18}weil auch Christus ein für alle  Mal für Sünden litt, \\
                    \spa{5}ein Gerechter für Ungerechte, \\
                    \spa{5}damit er euch zu Gott hinführe; \\
                    \spa{5}er wurde nämlich, \\
                    \spa{5}einerseits, \\
                        \spa{6}zu Tode gebracht am  Fleisch, \\
                    \spa{5}andererseits zum Leben gebracht \\
                        \spa{6}durch  den Geist, \\
                        \spa{6}\hh{19} in dem er auch hinging und den sich \\
                        \spa{6}in Gewahrsam befindenden \\
                        \spa{6}Geistern verkündete, \\
                        \spa{6}\hh{20}die im Unglauben ungehorsam waren, \\
                        \spa{6}einst, \\
                        \spa{6}als die Geduld Gottes am Warten war in den Tagen Noahs, \\
                        \spa{6}während eine Arche zubereitet wurde, \\
                        \spa{6}in die wenige, \\
                        \spa{6}das heißt, acht Seelen, \\
                        \spa{6}hinein[gingen und] hindurchgerettet wurden durch Wasser, \\
                    \spa{5}\hh{21}welches [als] Abbild nun auch uns  rettet und bewahrt, \\
                        \spa{6}[als] Taufe, \\
                        \spa{6}nicht ein Entfernen des Schmutzes am Fleisch, \\
                            \spa{7}sondern eine  verpflichtende Erklärung eines guten Gewissens, \\
                            \spa{7}an Gott [gerichtet], \\
                        \spa{6}– durch die Auferstehung Jesu Christi; \\
                            \spa{7}\hh{22}der ist zur Rechten Gottes, \\
                            \spa{7}nachdem er in den Himmel gegangen ist \\
                            \spa{7}und [himmlische] Boten \\
                            \spa{7}und Autoritäten \\
                            \spa{7}und Kräfte \\
                            \spa{7}ihm untergeordnet worden sind. \\
